\documentclass[../../main.tex]{subfiles}
% \documentclass[a4paper,10pt]{ltjsarticle}
\begin{document}
\rhead{京都大学後期$1994$年$6$番}
\setcounter{equation}{0}
\begin{oframed}
  \subfile{problem}
\end{oframed}
\setlength{\columnseprule}{0.4pt}
\begin{multicols}{2}
  {\bf [解答]}
  \paragraph{(1)}
  部分積分法より
  \begin{align}
    I_{n+1}
    & = \int_{1}^{e} (\log x)^{n+1} dx                                             \\
    & = \left[ x(\log x)^{n+1} \right]_{1}^{e} - \int_{1}^{e} (n+1)(\log x)^{n} dx \\
    & = e - (n+1)I_n
  \end{align}
  だから，求める漸化式は
  \begin{align}
    I_{n+1} =  e - (n+1)I_n \label{eq:1}
  \end{align}
  である．$\cdots$(答)

  \vspace{10pt}
  \paragraph{(2)}
  まずは下側の不等式から評価する．
  $[1,e]$ で $0 \le \log x \le 1$ だから,
  \begin{align}
    (\log x)^n \ge (\log x)^{n+1}
  \end{align}
  より，同じ区間で積分して $I_n \ge I_{n+1}$である．
  これに (1) を代入して
  \begin{align}
    I_n \ge e - (n+1)I_n \\
    \therefore I_n \ge \dfrac{e}{n+2} \quad (\because n \in \mathbb{N})
  \end{align}
  であり，ここで
  \begin{align}
    & \dfrac{e}{n+2} \ge \dfrac{e-1}{n+1} \\
    \iff & e(n+1) \ge (e-1)(n+2)             \\
    \iff & n+2     \ge e
  \end{align}
  で，$n\ge 1$よりこれは成立するから，
  \begin{align}
    & I_n \ge \dfrac{e}{n+2} \ge \dfrac{e-1}{n+1} \\
    \therefore
    & I_n \ge \dfrac{e-1}{n+1}
  \end{align}
  となり，下側の不等式の成立が示された．

  次に上側の評価を考える．漸化式\cref{eq:1}を二回利用して
  \begin{align}
    I_{n+2}
    & =  e - (n+2)I_{n+1}                 \\
    & =  e - (n+2)\left[e-(n+1)I_n\right] \\
    & =  -e(n+1) + (n+1)(n+2)I_n
  \end{align}
  であり，ここで$I_{n+2}\le I_{1} =1$だから
  \begin{align}
    & -e(n+1) + (n+1)(n+2)I_n \le 1       \\
    \therefore
    & I_n \le \dfrac{e(n+1)+1}{(n+1)(n+2)}
  \end{align}
  である．よって上側の不等式は示された．

  以上より，上下の不等式が成立することが示された．$\cdots$(答)

  \vspace{10pt}
  {\bf [別解]}

  \subparagraph{(2)}
  少し長くなるが数学的帰納法で示す．すなわち
  \begin{align}
    \dfrac{e-1}{n+1} \le I_n \le \dfrac{(n+1)e+1}{(n+1)(n+2)} \label{eq:2}
  \end{align}
  が全ての自然数$n$で成立することを示す．以下の補題を用意する．

  \begin{lemma}
    $\dfrac{e}{n+2} \le I_n$
  \end{lemma}

  \begin{proof}
    $[1,e]$で$0 \le \log x \le 1$だから，$(\log x)^n \ge (\log x)^{n+1}$より，
    同じ区間で積分して$I_n \ge I_{n+1}$である．これに(1)を代入して
    \begin{align}
      I_n \ge e-(n+1)I_n \\
      \therefore I_n \ge \dfrac{e}{n+2} \quad (\because n \in \mathbb{N})
    \end{align}
    となり，示された．
  \end{proof}

  \begin{lemma}
    $\dfrac{e}{n+2} \ge \dfrac{(n+1)^2e-1}{(n+1)^2(n+2)}$
  \end{lemma}

  \begin{proof}
    \begin{align}
      (\text{左辺})-(\text{右辺}) = \dfrac{1}{(n+1)^2(n+2)} > 0 \quad (\because n \in \mathbb{N})
    \end{align}
    から示された．
  \end{proof}

  \textbf{1$^\circ$ $n=1$の時}

  \begin{align}
    I_1 = \left[ x(\log x - 1) \right]_{1}^{e} = 1
  \end{align}
  であり，$2.7<e<2.8$から
  \begin{align}
    \begin{cases}
      1 - \dfrac{e-1}{2}          > 1 - \dfrac{1.8}{2} > 0 \\
      \dfrac{2e+1}{2 \cdot 3} - 1 > \dfrac{5.4-5}{6} > 0
    \end{cases}
  \end{align}
  となり，$n=1$で\cref{eq:2}は成立する．

  \textbf{2$^\circ$ $n=k \ (k \in \mathbb{N})$での成立を仮定する時}

  $n=k+1$での\cref{eq:2}の成立，つまり
  \begin{align}
    \dfrac{e-1}{k+2} \le I_{k+1} \le \dfrac{(k+1)e+1}{(k+1)(k+2)}
  \end{align}
  を示す．\cref{eq:1}より$I_{k+1}=e-(k+1)I_k$だから，これは
  \begin{align}
    \dfrac{e-1}{k+2} \le e-(k+1)I_k \le \dfrac{(k+1)e+1}{(k+1)(k+2)}
  \end{align}
  と同値であり，$k \in \mathbb{N}$より$k+1>0$に注意して$I_k$について解くと
  \begin{align}
    \dfrac{(k+1)^2e-1}{(k+1)^2(k+2)} \le I_k \le \dfrac{(k+1)e+1}{(k+1)(k+2)} \label{eq:3}
  \end{align}
  を示せば良い．\cref{eq:3}の右側の不等式は仮定から成立するので，左側の不等式を示す．

  [補題1]，[補題2]で$n=k$とすると，\cref{eq:3}の左側の不等式が示された．
  従って，$n=k+1$でも\cref{eq:2}は成立する．

  以上1$^\circ$，2$^\circ$から，全ての自然数$n$で\cref{eq:2}が成立する．従って題意は示された．$\cdots$(答)

  \vspace{10pt}
  {\bf [解説]}

  \newpage
\end{multicols}
\end{document}
